
% -------------------------------------------------------
%  Abstract
% -------------------------------------------------------


\begin{center}
\textbf{چکیده}
\end{center}

\begingroup
\setlength{\parskip}{0.22em}
\linespread{1.45}\selectfont
\noindent
هدف این پروژه، طراحی، پیاده‌سازی و ارزیابی یک نمونه اولیه سامانه گفتار‌به‌گفتار
فارسی است که به‌صورت تمام‌دوطرفه کار می‌کند؛ یعنی میکروفون کاربر در هنگام پخش
پاسخ نیز فعال می‌ماند و سامانه می‌تواند بر پایه یک تصمیم آشکارساز کلاسیک، منبع
صوت در حال پخش را متوقف کند و گفتار جمع‌شده کاربر را به نوبت بعدی منتقل نماید.
برای کاهش ریسک اجرایی، دو مسیر به‌طور موازی بررسی شد: نخست تطبیق مستقیم مدل
متن‌باز گفتار‌به‌گفتار \en{Moshika} با داده فارسی، و دوم یک سامانه آبشاری شامل
بازشناسی گفتار فارسی با \en{NeMo}، تولید پاسخ متنی با
\en{Qwen3-4B-Instruct-2507} و تبدیل متن به گفتار با \en{Piper}.

از آرشیو داخلی محتوای گفتاری فارسی، یک مجموعه داده مکالمه مشتق‌شده شامل
\N{6,754} جفت «گفتار کاربر / پاسخ» با مدت نهایی \N{108.584} ساعت صوت استریو
ساخته شد. تفکیک آموزش، اعتبارسنجی و آزمون بر پایه شناسه نشست منبع انجام گرفت و
نشت گروهی مشاهده نشد. یک ماژول تشخیص وقفه نیز بر پایه ویژگی‌های انرژی، زیر‌و‌بمی
و \en{MFCC} به همراه یک طبقه‌بند تقویت گرادیانی پیاده‌سازی شد.

پس از عبور از یک مرحله صلاحیت هفت‌ردیفی ازپیش‌قفل‌شده، پاسخ‌گوی
\en{Qwen3-4B} با پرامپت \en{v2} در آزمون نهایی \N{40} ردیفی گروه‌ناهمپوشان از
هر هفت شرط خودکار ازپیش‌تعریف‌شده عبور کرد: هر \N{240} فراخوانی داور معتبر بود،
میانگین ارتباط پاسخ از \N{1.425} برای خط پایه \en{Qwen2.5-0.5B} به \N{3.150}
رسید، میانگین انسجام \N{3.325} شد و نرخ برد زوجی \N{67.5} درصد به‌دست آمد.
داور، مدل هم‌خانواده \en{Qwen3.8-27B-FP8} بود؛ بنابراین این سنجه جایگزین ارزیابی
انسانی نیست. زنجیره مکانیکی آبشاری با پاسخ‌گوی \en{Qwen2.5-0.5B} در پنل ثابت
اعتبارسنجی در هر \N{9} ردیف رونویسی فارسی، پاسخ فارسی و صوت غیرساکت تولید کرد
و از پاسخ جایگزین قاعده‌محور استفاده نکرد. بازبازشناسی همان \N{9} خروجی، نرخ
خطای نویسه \N{7.14} درصد و نرخ خطای واژه \N{30.41} درصد را نشان داد. این پنل
پس از ارتقای پاسخ‌گو به \en{Qwen3-4B} دوباره اجرا نشد.

در مقابل، مسیر مستقیم در شش نسخه آزمایشی، با آن‌که زیان آموزشی را به‌طور
پیوسته کاهش داد، هیچ چک‌پوینتی تولید نکرد که از شرط اجرای خودرگرسیو عبور کند؛
نسخه \en{v6.2} کاهش \N{32.20} درصدی زیان متنی درون‌آموزشی را نشان داد اما هر
چهار چک‌پوینت آن در آزمون نه‌ردیفی اجرا \N{0} از \N{9} گرفتند. این مسیر
به‌عنوان یک نتیجه منفی با سیگنال یادگیری مثبت گزارش شده است و سامانه آبشاری
نمونه اولیه نهایی پروژه در نظر گرفته شد. آزمون مرورگر روی پردازنده گرافیکی
فیزیکی \N{12} تا \N{24} گیگابایتی، اثبات تأخیر سرتاسری حداکثر
\N{500} میلی‌ثانیه، برچسب‌گذاری مستقل انسانی وقفه و مطالعه \N{5} تا
\N{10} کاربر در مهلت پروژه انجام نشد و به‌عنوان محدودیت صریح باقی می‌ماند.

\vspace{0.6em}
\noindent \textbf{کلیدواژه‌ها}:
گفتار‌به‌گفتار، تمام‌دوطرفه، تشخیص وقفه، زبان فارسی، مدل زبانی بزرگ،
سامانه آبشاری، تطبیق کم‌پارامتر
\endgroup
\newpage
